\section{Method}
\label{sec:method}
\input{images_tex/pipeline}

\subsection{Overview}
\label{sec:method:overview}

We propose \ourMethod{}, a feed-forward framework for pairwise 3D
reconstruction from low-light RGB images and temporally aligned event streams.
The framework preserves the encoder--decoder interface of a pretrained
MASt3R model~\cite{dust3r_cvpr24,mast3r_eccv24} and introduces a separate event
encoder together with a Bounded Reliability Gate (BRG).  The RGB encoder
remains the primary geometry pathway.  Event features are injected once, after
the two per-view encoders and before the pretrained cross-view decoder, as
small token-wise residuals.

For each view, the event encoder maps an eight-channel polarity-separated
voxel grid to the same spatial token grid and feature dimension as the RGB
encoder.  BRG derives a scalar gate for every token from local event activity,
occupancy, and polarity statistics.  A scheduled gate cap and a bounded global
event scale constrain the coefficient of each event residual.  The fused token
pair is then processed by the original cross-view decoder and prediction heads,
which produce dense point maps, point confidence, local descriptors, and
descriptor confidence.

During training, a frozen RGB-only model provides descriptor and point-map
consistency targets.  The teacher receives the stored source RGB pair before
online low-light degradation, whereas the event student receives the degraded
pair.  For native nighttime data on which additional degradation is disabled,
the two branches receive the same RGB pair and differ only in the event path.
Metric depth, depth-gradient, event-edge, and event-photometric terms are
enabled according to the metadata available in each source.  Their active
weights are normalized to keep the objective scale comparable across sources.
At inference time, the teacher and all losses are removed.

The overall architecture is illustrated in \cref{fig:pipeline}.


\subsection{Problem Formulation and Input Construction}
\label{sec:method:formulation}

\noindent\textbf{Pairwise point-map prediction.}
Let
$\mathcal{I}^{\ell}=(I_1^{\ell},I_2^{\ell})$ be a low-light RGB pair and
$\mathcal{V}=(\mathbf{V}_1,\mathbf{V}_2)$ the corresponding event voxels.
The student predicts
\begin{equation}
\mathcal{O}_{S}=f_{\theta}(\mathcal{I}^{\ell},\mathcal{V})
=\left\{
\hat{\mathbf{X}}_{S}^{v,1},
\hat{\mathbf{C}}_{S}^{v,1},
\hat{\mathbf{D}}_{S}^{v},
\hat{\mathbf{Q}}_{S}^{v}
\right\}_{v\in\{1,2\}}.
\label{eq:student_outputs}
\end{equation}
Following DUSt3R and MASt3R, both point maps are expressed in the coordinate
frame of the first camera.  Thus,
$\hat{\mathbf{X}}_{S}^{1,1}$ is the point map of the first image and
$\hat{\mathbf{X}}_{S}^{2,1}$, implemented as
\texttt{pts3d\_in\_other\_view}, is the point map of the second image expressed
in the first-view frame.  The remaining outputs are point confidence
$\hat{\mathbf{C}}$, dense matching descriptor $\hat{\mathbf{D}}$, and
descriptor confidence $\hat{\mathbf{Q}}$.

\vspace{+5pt}
\noindent\textbf{Event voxelization.}
An event is represented as
$e_k=(x_k,y_k,t_k,p_k)$, where $p_k\in\{-1,+1\}$.  For the events selected in
one temporal window, we normalize timestamps using the first and last event in
that window,
\begin{equation}
t_k^{*}
=(B-1)\frac{t_k-t_{\min}}{t_{\max}-t_{\min}+\epsilon}.
\label{eq:event_time_normalization}
\end{equation}
With $\kappa(a)=\max(0,1-|a|)$, temporal bilinear accumulation gives
\begin{equation}
\mathbf{V}^{p}(b,x,y)
=\sum_{k:p_k=p}
\delta_{x,x_k}\delta_{y,y_k}\kappa(b-t_k^{*}),
\qquad b\in\{0,\ldots,B-1\}.
\label{eq:event_voxelization}
\end{equation}
We use $B=4$ bins per polarity and stack the positive bins before the negative
bins, yielding
$\mathbf{V}\in\mathbb{R}^{8\times H\times W}$.  Before entering the event
encoder, the adapter applies signed logarithmic compression and normalizes only
nonzero entries,
\begin{equation}
\overline{\mathbf{V}}
=\operatorname{NZNorm}\!\left(
\operatorname{sign}(\mathbf{V})\log(1+|\mathbf{V}|)
\right),
\label{eq:event_voxel_normalization}
\end{equation}
so empty pixels remain zero.  Local gate statistics, defined below, are
computed from the input voxel before this adapter-side compression.

Native-event datasets are converted to the same per-frame voxel format before
training.  For DL3DV, which does not contain native events, we first densify the
RGB sequence with P-RIFE and then simulate the event stream using
DVS-Voltmeter.  The simulated stream is partitioned into frame-associated
windows and voxelized with the same eight-channel contract.  This simulation
is an offline data-preparation step and is not used at inference time.

\vspace{+5pt}
\noindent\textbf{Online low-light degradation.}
For sources on which synthetic darkening is enabled, the stored RGB image
$I^{o}\in[0,1]$ is mapped to an approximate linear domain and degraded as
\begin{align}
I^{\mathrm{lin}} &= (I^{o})^{\gamma}, \qquad \gamma=2.2, \\
\widetilde I^{\mathrm{lin}}
&=eI^{\mathrm{lin}}
+\sigma_s\sqrt{eI^{\mathrm{lin}}+\epsilon}\,\xi_s
+\sigma_r\xi_r, \\
I^{\ell}
&=\operatorname{clip}(\widetilde I^{\mathrm{lin}},0,1)^{1/\gamma},
\label{eq:lowlight_degradation}
\end{align}
where $\log e$ is sampled uniformly between the configured exposure bounds and
$\xi_s,\xi_r\sim\mathcal{N}(0,1)$.  This is an approximate linear-RGB exposure
model with signal-dependent Gaussian shot noise and additive Gaussian read
noise; it is not a claim of processing true camera RAW data.  For real
nighttime adaptation, this augmentation can be disabled, in which case
$I^{\ell}=I^{o}$.


\subsection{Reliability-Aware Event Residual Adaptation}
\label{sec:method:event_integration}

For clarity, this subsection omits the view index; the same per-view operation
is applied to both inputs.

\noindent\textbf{Event token encoder.}
The pretrained image encoder produces
\begin{equation}
\mathbf{Z}^{I}=\Phi_I(I^{\ell})
\in\mathbb{R}^{N\times d}.
\label{eq:image_tokens}
\end{equation}
The event voxel is patchified by a convolution whose kernel and stride match
the image patch size.  The resulting tokens are processed by a separate copy
of the pretrained image-encoder blocks,
\begin{equation}
\mathbf{Z}^{E}=\Phi_E(\overline{\mathbf{V}})
\in\mathbb{R}^{N\times d}.
\label{eq:event_tokens}
\end{equation}
The copied blocks are trainable, but their architecture and positional token
grid match those of the RGB encoder.  No event-to-image reconstruction is
performed.

\vspace{+5pt}
\noindent\textbf{Local event statistics.}
\input{images_tex/brg_gate}
Let $\Omega_i$ be the pixels represented by token $i$.  At pixel
$\mathbf{p}=(x,y)$, define
\begin{align}
a(\mathbf{p})
&=\sum_{c=1}^{2B}|\mathbf{V}_{c}(\mathbf{p})|, \\
a^{+}(\mathbf{p})
&=\sum_{b=0}^{B-1}|\mathbf{V}^{+}_{b}(\mathbf{p})|, \\
a^{-}(\mathbf{p})
&=\sum_{b=0}^{B-1}|\mathbf{V}^{-}_{b}(\mathbf{p})|.
\end{align}
We compute three pixel-wise maps and average-pool them to the event-token grid:
\begin{equation}
\mathbf{r}_{i}=
\begin{bmatrix}
\displaystyle
\frac{1}{|\Omega_i|}
\sum_{\mathbf{p}\in\Omega_i}\log(1+a(\mathbf{p})) \\[8pt]
\displaystyle
\frac{1}{|\Omega_i|}
\sum_{\mathbf{p}\in\Omega_i}\mathbb{I}[a(\mathbf{p})>0] \\[8pt]
\displaystyle
\frac{1}{|\Omega_i|}
\sum_{\mathbf{p}\in\Omega_i}
\frac{|a^{+}(\mathbf{p})-a^{-}(\mathbf{p})|}
{a^{+}(\mathbf{p})+a^{-}(\mathbf{p})+\epsilon}
\end{bmatrix}.
\label{eq:reliability_statistics}
\end{equation}
The entries are log activity, occupancy, and polarity imbalance.  In
particular, activity and polarity are computed per pixel before pooling; they
are not obtained by first summing the whole token region.  We normalize the
statistics as
\begin{equation}
\hat{\mathbf{r}}_i
=\boldsymbol{\gamma}\odot
\operatorname{clip}\!\left(
\frac{\mathbf{r}_i-\boldsymbol{\mu}}
{\boldsymbol{\sigma}+10^{-4}},-K,K
\right)+\boldsymbol{\beta},
\label{eq:event_stat_normalization}
\end{equation}
where $\boldsymbol{\mu}$ and $\boldsymbol{\sigma}$ are configured statistic
moments.  The affine parameters $\boldsymbol{\gamma},\boldsymbol{\beta}$ are
learned.  Dataset-specific event-rate differences are handled by the learned
source gate bias in \cref{eq:gate_logit}, rather than by changing the event
content branch.

\vspace{+5pt}
\noindent\textbf{Bounded Reliability Gate.}
The event content and its reliability coefficient are computed separately:
\begin{align}
\Delta\mathbf{Z}^{E}
&=\phi_r\!\left(\operatorname{LN}(\mathbf{Z}^{E})\right),
\label{eq:event_residual}\\
\ell_i
&=\phi_g(\hat{\mathbf{r}}_i)+\beta_s,
\label{eq:gate_logit}
\end{align}
where $\phi_r$ is a two-layer MLP, $\phi_g$ is a lightweight gate MLP, and
$\beta_s$ is a learned source bias when source conditioning is enabled.  The
raw logit is smoothly bounded and converted to a scheduled token gate,
\begin{align}
\bar\ell_i
&=L\tanh(\ell_i/L), \\
g_i(u)
&=c(u)\,\sigma(\bar\ell_i/\tau)\,
\mathbb{I}[o_i>\epsilon_o],
\label{eq:bounded_reliability_gate}
\end{align}
where $o_i$ is the occupancy entry of $\mathbf{r}_i$, $c(u)$ is linearly
scheduled from an initial to a final gate cap, and $\tau$ is the gate
temperature.  A second sigmoid parameterization bounds the global event scale,
\begin{equation}
s_e=s_{\max}\sigma(a).
\label{eq:event_global_scale}
\end{equation}
With channel modulation disabled in the reported model, fusion is
\begin{equation}
\widetilde{\mathbf{Z}}^{I}_{i}
=\mathbf{Z}^{I}_{i}
+s_e g_i(u)\Delta\mathbf{Z}^{E}_{i}.
\label{eq:gated_fusion}
\end{equation}
Therefore the scalar coefficient multiplying each event residual is bounded by
$s_{\max}c(u)$.  The bound constrains the residual coefficient, not the norm of
$\Delta\mathbf{Z}^{E}_{i}$ itself.

BRG contains no RGB--event cross-attention.  The residual
$\Delta\mathbf{Z}^{E}$ depends on the event token, while its gate depends on
event statistics.  Cross-view interactions occur only after fusion, when the
frozen MASt3R decoder processes the two adapted view-token sequences,
\begin{equation}
\mathcal{O}_{S}
=\Psi\!\left(
\widetilde{\mathbf{Z}}^{I}_{1},
\widetilde{\mathbf{Z}}^{I}_{2}
\right).
\label{eq:adapted_prediction}
\end{equation}
The scalar gated cross-attention module is retained only as an experimental
baseline and is not part of the final BRG architecture.


\subsection{Frozen RGB Geometry Guidance}
\label{sec:method:distillation}

Let $\mathcal{I}^{o}=(I_1^{o},I_2^{o})$ denote the stored source RGB pair before
online low-light degradation.  A frozen RGB-only teacher predicts
\begin{equation}
\mathcal{O}_{T}
=f_{\theta_T}(\mathcal{I}^{o})
=\left\{
\hat{\mathbf{X}}_{T}^{v,1},
\hat{\mathbf{C}}_{T}^{v,1},
\hat{\mathbf{D}}_{T}^{v},
\hat{\mathbf{Q}}_{T}^{v}
\right\}_{v\in\{1,2\}}.
\label{eq:teacher_prediction}
\end{equation}
The teacher pair and student pair correspond pixel-wise because the latter is
obtained from the former by photometric degradation.  The source pair is
well-exposed for synthetic low-light sources but is not assumed to be
well-exposed for native nighttime datasets.  Teacher parameters and outputs
are detached throughout training.

\vspace{+5pt}
\noindent\textbf{Dense descriptor distillation.}
We align the final dense descriptor maps returned by the MASt3R heads rather
than collecting intermediate decoder activations.  Let
$\mathcal{P}^{D}_{v}$ be locations where both student and teacher descriptors
are finite.  The descriptor loss is
\begin{equation}
\mathcal{L}_{\mathrm{desc}}
=\frac{1}{2}\sum_{v=1}^{2}
\frac{1}{|\mathcal{P}^{D}_{v}|}
\sum_{\mathbf{p}\in\mathcal{P}^{D}_{v}}
\left[
1-\operatorname{cos}\!\left(
\hat{\mathbf{D}}_{S}^{v}(\mathbf{p}),
\operatorname{sg}(\hat{\mathbf{D}}_{T}^{v}(\mathbf{p}))
\right)
\right].
\label{eq:descriptor_distillation}
\end{equation}

\vspace{+5pt}
\noindent\textbf{Confidence-masked point-map consistency.}
For view $v$, we retain finite teacher/student points whose teacher confidence
exceeds $\tau_c$ and whose point norms remain below a fixed outlier threshold.
The valid set is denoted $\mathcal{P}^{X}_{v}$.  A robust scale is estimated
\emph{separately for each view},
\begin{equation}
s_{T,v}
=\operatorname{clip}\!\left(
\operatorname*{median}_{\mathbf{p}\in\mathcal{P}^{X}_{v}}
\left\|
\hat{\mathbf{X}}_{T}^{v,1}(\mathbf{p})
\right\|_2,
s_{\min},s_{\max}
\right).
\label{eq:teacher_geometry_scale}
\end{equation}
For coordinate $c\in\{1,2,3\}$, define the clipped normalized residual
\begin{equation}
r_{v,\mathbf{p},c}
=\operatorname{clip}\!\left(
\frac{
[\hat{\mathbf{X}}_{S}^{v,1}(\mathbf{p})
-\operatorname{sg}(\hat{\mathbf{X}}_{T}^{v,1}(\mathbf{p}))]_c
}{s_{T,v}},
-r_{\max},r_{\max}
\right).
\end{equation}
The implemented structure-consistency penalty is
\begin{equation}
\mathcal{L}_{\mathrm{struct}}
=\frac{1}{|\mathcal{V}_{X}|}
\sum_{v\in\mathcal{V}_{X}}
\underset{\substack{
\mathbf{p}\in\mathcal{P}^{X}_{v}\\c\in\{1,2,3\}
}}{\operatorname{mean}}
\log(1+|r_{v,\mathbf{p},c}|),
\label{eq:pointmap_consistency_loss}
\end{equation}
where $\mathcal{V}_{X}$ contains views with enough valid support.  This is a
bounded-influence consistency target rather than metric ground truth.  We use
the $\log(1+|r|)$ penalty to suppress the influence of severe point outliers.


\subsection{Source-Aware Training Objective}
\label{sec:method:training_obj}

\noindent\textbf{Active-weight normalization.}
Let
\begin{equation}
\mathcal{K}
=\{\mathrm{desc},\mathrm{struct},\mathrm{edge},\mathrm{photo},
\mathrm{depth},\mathrm{grad},\mathrm{gate}\}
\end{equation}
be the training objectives.  For a batch $\mathcal{B}$ from source $s$,
the effective component weight is
\begin{equation}
w_k(\mathcal{B})
=a_k(\mathcal{B})\,\alpha_{s,k}\lambda_k,
\label{eq:effective_loss_weight}
\end{equation}
where $a_k$ records whether the required tensors and valid pixels are present,
$\alpha_{s,k}$ is the configured source factor, and $\lambda_k$ is the base
weight.  The loss is normalized by the sum of weights of components that were
actually evaluated,
\begin{equation}
\mathcal{L}_{\mathrm{train}}(\mathcal{B})
=\frac{\lambda_{\mathrm{target}}}
{\sum_{k\in\mathcal{K}}w_k(\mathcal{B})+\epsilon}
\sum_{k\in\mathcal{K}}
w_k(\mathcal{B})\mathcal{L}_k.
\label{eq:dynamic_loss}
\end{equation}
In the default configuration $\lambda_{\mathrm{target}}=1$.  Source factors
set unsupported terms to zero; for example, DL3DV does not receive metric-depth
supervision, while event-photometric supervision is used only for sources with
valid intrinsics and relative pose.

\vspace{+5pt}
\noindent\textbf{Metric depth supervision.}
We apply metric depth supervision to the first-view prediction using the $z$
coordinate of
$\hat{\mathbf{X}}_{S}^{1,1}$,
\begin{equation}
\hat D_1(\mathbf{p})
=\left[\hat{\mathbf{X}}_{S}^{1,1}(\mathbf{p})\right]_z,
\end{equation}
and valid measured depth $D_1^{\mathrm{gt}}$.  With
$d(\mathbf{p})=\log\hat D_1(\mathbf{p})-\log D_1^{\mathrm{gt}}(\mathbf{p})$,
the implemented SILog loss is
\begin{equation}
\mathcal{L}_{\mathrm{depth}}
=10\sqrt{
\operatorname{Var}_{\mathcal{P}^{D}}[d]
+\beta\left(
\operatorname{Mean}_{\mathcal{P}^{D}}[d]
\right)^2
+\epsilon},
\qquad \beta=0.15.
\label{eq:metric_depth_loss}
\end{equation}
This term directly anchors the metric scale of the first-view point map; the
second-view point map is instead constrained by teacher consistency and the
event-edge objective.

The local gradient term compares predicted and measured log-depth gradients.
Let
\begin{equation}
E(\mathbf{p})
=\omega_I\|\nabla I_1^{\ell}(\mathbf{p})\|_1
+\omega_E\|\nabla A_1(\mathbf{p})\|_1,
\qquad
A_v=\sum_c|\mathbf{V}_{v,c}|,
\end{equation}
and let $W=1+[\operatorname{ZNorm}(E)]_+$.  We use a validity-masked,
edge-weighted $L_1$ difference along horizontal and vertical directions,
\begin{equation}
\mathcal{L}_{\mathrm{grad}}
=\frac{1}{2}
\sum_{q\in\{x,y\}}
\operatorname{WMean}_{\mathcal{P}^{D}_{q}}
\left(
\left|
\nabla_q\log\hat D_1-
\nabla_q\log D_1^{\mathrm{gt}}
\right|;W_q
\right).
\label{eq:depth_gradient_loss}
\end{equation}

\vspace{+5pt}
\noindent\textbf{Event--geometry edge consistency.}
For both output point maps, let $\hat Z_v$ denote their $z$-coordinate map in
the common first-view frame.  The event-edge term correlates event activity
with the magnitude of the corresponding geometry gradient,
\begin{equation}
\mathcal{L}_{\mathrm{edge}}
=\frac{1}{2}\sum_{v=1}^{2}
\left[
1-\operatorname{corr}\!\left(
A_v,\|\nabla\hat Z_v\|_1
\right)
\right].
\label{eq:event_edge_loss}
\end{equation}
Both maps are z-normalized over finite pixels before evaluating the
correlation.  This is weak structural guidance; it does not assume that every
event is a depth discontinuity.

\vspace{+5pt}
\noindent\textbf{Event-photometric consistency.}
When camera intrinsics and a relative pose are available, the first-view
predicted depth warps the second log-grayscale image to the first view.  Define
\begin{align}
\Delta L^{I}
&=\mathcal{W}(\log \bar I_2^{\ell};\hat D_1,K,T_{1\rightarrow2})
-\log \bar I_1^{\ell}, \\
\Delta L^{E}
&=\sum_{b=0}^{B-1}\mathbf{V}^{+}_{1,b}
-\sum_{b=0}^{B-1}\mathbf{V}^{-}_{1,b},
\end{align}
where $\bar I$ is the channel-mean intensity.  On valid warped pixels, both
changes are standardized and compared by cosine similarity,
\begin{equation}
\mathcal{L}_{\mathrm{photo}}
=1-\operatorname{cos}\!\left(
\operatorname{ZNorm}(\Delta L^{I}),
\operatorname{ZNorm}(\Delta L^{E})
\right).
\label{eq:event_photometric_loss}
\end{equation}
This comparison is invariant to affine scale in the two change maps and is
therefore used as a weak consistency signal rather than a calibrated
pixel-wise brightness reconstruction objective.

\vspace{+5pt}
\noindent\textbf{Gate regularization.}
Let $\bar g(u)$ be the scheduled cap on the mean gate and let
$\mathcal{P}_{\mathrm{low}}=\{i:o_i<o_{\mathrm{thr}}\}$.  The reported model
uses
\begin{equation}
\begin{aligned}
\mathcal{L}_{\mathrm{gate}}
={}&\mu_{\mathrm{mean}}
\left[\operatorname{mean}_i g_i-\bar g(u)\right]_{+}^{2}\\
&+\mu_{\mathrm{low}}
\operatorname{mean}_{i\in\mathcal{P}_{\mathrm{low}}}g_i^2
+\mu_{\mathrm{logit}}
\operatorname{mean}_i
\left[|\ell_i|-\ell_{\mathrm{reg}}\right]_{+}^{2}.
\end{aligned}
\label{eq:gate_regularization}
\end{equation}
The terms discourage a globally open gate, suppress residuals in weak-event
regions, and prevent raw gate logits from saturating.  Channel modulation and
RGB-visibility gating are disabled in the reported BRG configuration.

\vspace{+5pt}
\noindent\textbf{Optimization and inference.}
The RGB image encoder, MASt3R cross-view decoder, prediction heads, and teacher
are frozen.  We optimize the copied event encoder, event residual projection,
gate network, global event scale, learned statistic affine, and source gate
biases.  Gradients pass through the frozen cross-view decoder to the event
adapter.  At inference time, only the low-light RGB pair, event voxels, RGB
branch, event encoder, BRG, and original decoder/heads are retained.
